\section{Conclusion and Future Work}
\label{sec:conclusion}

\subsection{Summary}
\label{subsec:concl-summary}
This thesis studied the Sim2Real gap for vision-based runway detection with a
YOLOv8-pose model fine-tuned on nominal-condition synthetic data from LARD-V2 and
evaluated on a weather-stratified real test set from LARD-V1. Under a single,
lightweight, reproducible protocol (fixed 5000-image budget, shared filtered test set,
primary metric mAP@[.50:.95]), it first \emph{quantified} the gap and then compared
three strategies to \emph{close} it.

The headline findings are: (RQ1) once the test sets are matched on both runway size and
\emph{vertical angle}, a modest but genuine Sim2Real gap appears --- previously hidden
because the real set looked deceptively easy --- and much of the residual synthetic
difficulty is an image-quality effect concentrated on distant runways; the combined
source generalizes best overall, with different sources suiting different weather
scenarios and fog the hardest. (RQ2) Weather augmentation gives limited and unstable
gains, the scenario-matched augmentation is often not the best, and it does not
substitute for source diversity. (RQ3) Self-training on unlabeled real landing sequences
is the most effective technique so far, improving adverse LARD splits by over $10$ mAP
and in-distribution sequences by up to $35$ mAP, though it is sensitive to pseudo-label
noise. (RQ4) Control-guided diffusion style transfer, designed to import real appearance
while preserving the corner labels, is ongoing.

\subsection{Contributions Revisited}
\label{subsec:concl-contrib}
Against the LARD-V2 study, which analyzed cross-dataset generalization on synthetic test
sets, this thesis contributes: a real-domain re-analysis of source suitability; a
weather-stratified real benchmark extended with private landing sequences to compensate
for sparse adverse-weather data; the vertical-angle filtering needed to measure the gap
fairly; and a head-to-head comparison of weather augmentation, self-training and
style transfer for lightweight runway-corner detection.

\subsection{Future Work}
\label{subsec:concl-future}
Several threads remain open. The \textbf{tracker-augmented self-training} variant
(\Cref{subsec:pl-tracker}) should be completed and compared against the confidence-only
loop, quantifying the benefit of temporal consistency for distant runways. The
\textbf{diffusion style-transfer} stage should be evaluated and, time permitting,
contrasted with a CycleGAN baseline to measure both accuracy and label integrity.
Beyond these, promising directions include \emph{combining} the techniques (e.g.\
self-training on style-transferred data), extending the evaluation to full pose/geometry
metrics, enlarging and balancing the adverse-weather real splits, and closing the
far-runway fidelity gap identified in \Cref{subsec:gap-quality}.
